% !TEX program = pdflatex
% Анализ, верификация и сопоставление реализации непараметрической
% знаково-ранговой карты Шухарта (SR) со средствами языка R.
% Первоисточники: Bakir (2004); Gibbons & Chakraborti (2010);
% Chakraborti & Graham (2019); Nascimento et al. (2026).
% Сборка: pdflatex (T2A + babel russian).
\documentclass[11pt,a4paper]{article}

\usepackage[T2A]{fontenc}
\usepackage[utf8]{inputenc}
\usepackage[russian]{babel}
\usepackage[margin=2.4cm]{geometry}
\usepackage{amsmath,amssymb}
\usepackage{pgfplots}
\pgfplotsset{compat=1.17}
\usepackage{booktabs}
\usepackage{array}
\usepackage{graphicx}
\usepackage{xcolor}
\usepackage{listings}
\usepackage{caption}
\usepackage{enumitem}
\usepackage[hidelinks]{hyperref}

\definecolor{accent}{RGB}{46,134,171}   % ACCENT  #2E86AB
\definecolor{alert}{RGB}{178,34,34}     % ALERT   #B22222
\definecolor{third}{RGB}{60,110,71}     % PALETTE #3C6E47
\definecolor{softgray}{RGB}{128,128,128} % CENTER #808080

\lstset{
  basicstyle=\ttfamily\small,
  commentstyle=\itshape\color{softgray},
  keywordstyle=\bfseries,
  breaklines=true,
  frame=single,
  rulecolor=\color{softgray},
  numbers=left,
  numberstyle=\tiny\color{softgray},
  numbersep=8pt,
  columns=fullflexible,
  keepspaces=true,
  upquote=true,
}

\captionsetup{font=small,labelfont=bf}
\setlist[itemize]{itemsep=1pt,topsep=2pt}
\setlist[description]{itemsep=1pt,topsep=2pt}
\emergencystretch=1.5em

\newcommand{\SR}{\mathrm{SR}}
\newcommand{\FAR}{\mathrm{FAR}}
\newcommand{\ARL}{\mathrm{ARL}}
\newcommand{\MRL}{\mathrm{MRL}}
\newcommand{\FARz}{\mathrm{FAR}_0}
\newcommand{\SD}{\operatorname{SD}}
\newcommand{\far}{\mathrm{far}}

\title{\bfseries Непараметрическая знаково-ранговая карта Шухарта (SR):\\
математические основы, программная реализация,\\
верификация и сопоставление с языком R}
\author{Проект \texttt{nonparametric-charts}}
\date{\today}

% ==================== Координаты графиков (sr_gen_figdata.py) ====================
% ==== SR: границы a(n) по режимам, n = 2..20 ====
\def\vSRACons{(2,3.000000) (3,6.000000) (4,10.000000) (5,15.000000) (6,21.000000) (7,28.000000) (8,36.000000) (9,45.000000) (10,55.000000) (11,64.000000) (12,72.000000) (13,81.000000) (14,91.000000) (15,100.000000) (16,110.000000) (17,121.000000) (18,133.000000) (19,144.000000) (20,156.000000)}
\def\vSRASym{(2,3.000000) (3,6.000000) (4,10.000000) (5,15.000000) (6,21.000000) (7,28.000000) (8,36.000000) (9,45.000000) (10,55.000000) (11,66.000000) (12,78.000000) (13,87.000000) (14,97.000000) (15,106.000000) (16,118.000000) (17,127.000000) (18,139.000000) (19,150.000000) (20,162.000000)}
\def\vSRANear{(2,2.000000) (3,5.000000) (4,9.000000) (5,14.000000) (6,18.000000) (7,23.000000) (8,27.000000) (9,32.000000) (10,38.000000) (11,43.000000) (12,49.000000) (13,56.000000) (14,62.000000) (15,69.000000) (16,75.000000) (17,82.000000) (18,90.000000) (19,97.000000) (20,105.000000)}
\def\vSRANpExact{(2,3.000000) (3,6.000000) (4,10.000000) (5,15.000000) (6,21.000000) (7,28.000000) (8,36.000000) (9,45.000000) (10,53.000000) (11,62.000000) (12,70.000000) (13,79.000000) (14,89.000000) (15,98.000000) (16,108.000000) (17,119.000000) (18,131.000000) (19,142.000000) (20,154.000000)}
\def\vSRANpAsym{(2,6.708152) (3,11.224886) (4,16.431551) (5,22.248425) (6,28.617957) (7,35.496206) (8,42.848242) (9,50.645441) (10,58.863799) (11,67.482814) (12,76.484706) (13,85.853870) (14,95.576461) (15,105.640091) (16,116.033588) (17,126.746809) (18,137.770493) (19,149.096140) (20,160.715910)}
% ==== SR: достигнутый FAR(n) по режимам, n = 2..20 ====
\def\vSRFARCons{(2,0.500000) (3,0.250000) (4,0.125000) (5,0.062500) (6,0.031250) (7,0.015625) (8,0.007812) (9,0.003906) (10,0.001953) (11,0.001953) (12,0.002441) (13,0.002441) (14,0.002319) (15,0.002625) (16,0.002686) (17,0.002579) (18,0.002335) (19,0.002399) (20,0.002325)}
\def\vSRFARSym{(2,0.500000) (3,0.250000) (4,0.125000) (5,0.062500) (6,0.031250) (7,0.015625) (8,0.007812) (9,0.003906) (10,0.001953) (11,0.000977) (12,0.000488) (13,0.000732) (14,0.000854) (15,0.001160) (16,0.001007) (17,0.001343) (18,0.001289) (19,0.001411) (20,0.001432)}
\def\vSRFARNear{(2,0.500000) (3,0.250000) (4,0.125000) (5,0.062500) (6,0.062500) (7,0.046875) (8,0.054688) (9,0.054688) (10,0.048828) (11,0.053711) (12,0.052246) (13,0.047852) (14,0.049438) (15,0.047913) (16,0.050659) (17,0.050537) (18,0.048279) (19,0.049366) (20,0.048441)}
% ==== SR: таблица режимов, n = 2..20 ====
\def\vSRModeTable{%
2 & 3 & 0.500000 & 3 & 0.500000 & 2 & 0.500000 \\
3 & 6 & 0.250000 & 6 & 0.250000 & 5 & 0.250000 \\
4 & 10 & 0.125000 & 10 & 0.125000 & 9 & 0.125000 \\
5 & 15 & 0.062500 & 15 & 0.062500 & 14 & 0.062500 \\
6 & 21 & 0.031250 & 21 & 0.031250 & 18 & 0.062500 \\
7 & 28 & 0.015625 & 28 & 0.015625 & 23 & 0.046875 \\
8 & 36 & 0.007812 & 36 & 0.007812 & 27 & 0.054688 \\
9 & 45 & 0.003906 & 45 & 0.003906 & 32 & 0.054688 \\
10 & 55 & 0.001953 & 55 & 0.001953 & 38 & 0.048828 \\
11 & 64 & 0.001953 & 66 & 0.000977 & 43 & 0.053711 \\
12 & 72 & 0.002441 & 78 & 0.000488 & 49 & 0.052246 \\
13 & 81 & 0.002441 & 87 & 0.000732 & 56 & 0.047852 \\
14 & 91 & 0.002319 & 97 & 0.000854 & 62 & 0.049438 \\
15 & 100 & 0.002625 & 106 & 0.001160 & 69 & 0.047913 \\
16 & 110 & 0.002686 & 118 & 0.001007 & 75 & 0.050659 \\
17 & 121 & 0.002579 & 127 & 0.001343 & 82 & 0.050537 \\
18 & 133 & 0.002335 & 139 & 0.001289 & 90 & 0.048279 \\
19 & 144 & 0.002399 & 150 & 0.001411 & 97 & 0.049366 \\
20 & 156 & 0.002325 & 162 & 0.001432 & 105 & 0.048441 \\}
% ==== SR: пример n = 14 (критерий nearest), a = 40..80 ====
\def\vSRNearAbs14{(40,0.216553) (41,0.216553) (42,0.193726) (43,0.193726) (44,0.172607) (45,0.172607) (46,0.153076) (47,0.153076) (48,0.135254) (49,0.135254) (50,0.118896) (51,0.118896) (52,0.104004) (53,0.104004) (54,0.090576) (55,0.090576) (56,0.078491) (57,0.078491) (58,0.067627) (59,0.067627) (60,0.057983) (61,0.057983) (62,0.049438) (63,0.049438) (64,0.041870) (65,0.041870) (66,0.035278) (67,0.035278) (68,0.029541) (69,0.029541) (70,0.024536) (71,0.024536) (72,0.020264) (73,0.020264) (74,0.016602) (75,0.016602) (76,0.013428) (77,0.013428) (78,0.010742) (79,0.010742) (80,0.008545)}
\def\vSRNearMin14{(62,0.049438) (63,0.049438)}
% ==== SR: мощность при сдвиге (n = 16), delta = 0..2 ====
% n=16: книга a=110, npSPC exact a=108
\def\vSROCBk{(0.0,0.002686) (0.25,0.014925) (0.5,0.086417) (0.75,0.295649) (1.0,0.623783) (1.25,0.882588) (1.5,0.980478) (1.75,0.998228) (2.0,0.999897)}
\def\vSROCNp{(0.0,0.003357) (0.25,0.017994) (0.5,0.100213) (0.75,0.328371) (1.0,0.662585) (1.25,0.904110) (1.5,0.985944) (1.75,0.998910) (2.0,0.999948)}
% ==== SR: распределение SR в IC (n = 5) ====
\def\vSRpmfFive{(-15,0.031250) (-13,0.031250) (-11,0.031250) (-9,0.062500) (-7,0.062500) (-5,0.093750) (-3,0.093750) (-1,0.093750) (1,0.093750) (3,0.093750) (5,0.093750) (7,0.062500) (9,0.062500) (11,0.031250) (13,0.031250) (15,0.031250)}
% ==== SR: односторонний хвост P(SR>=a) для n = 5,10,14,20 ====
\def\vSRTailFive{(15,0.031250) (13,0.062500) (11,0.093750) (9,0.156250) (7,0.218750) (5,0.312500) (3,0.406250) (1,0.500000)}
\def\vSRTailTen{(55,0.000977) (53,0.001953) (51,0.002930) (49,0.004883) (47,0.006836) (45,0.009766) (43,0.013672) (41,0.018555) (39,0.024414) (37,0.032227) (35,0.041992) (33,0.052734) (31,0.065430) (29,0.080078) (27,0.096680) (25,0.116211) (23,0.137695) (21,0.161133) (19,0.187500) (17,0.215820) (15,0.246094) (13,0.278320) (11,0.312500) (9,0.347656) (7,0.384766) (5,0.422852) (3,0.460938) (1,0.500000)}
\def\vSRTailFourteen{(105,0.000061) (103,0.000122) (101,0.000183) (99,0.000305) (97,0.000427) (95,0.000610) (93,0.000854) (91,0.001160) (89,0.001526) (87,0.002014) (85,0.002625) (83,0.003357) (81,0.004272) (79,0.005371) (77,0.006714) (75,0.008301) (73,0.010132) (71,0.012268) (69,0.014771) (67,0.017639) (65,0.020935) (63,0.024719) (61,0.028992) (59,0.033813) (57,0.039246) (55,0.045288) (53,0.052002) (51,0.059448) (49,0.067627) (47,0.076538) (45,0.086304) (43,0.096863) (41,0.108276) (39,0.120605) (37,0.133789) (35,0.147888) (33,0.162903) (31,0.178772) (29,0.195496) (27,0.213135) (25,0.231567) (23,0.250793) (21,0.270813) (19,0.291504) (17,0.312866) (15,0.334900) (13,0.357422) (11,0.380432) (9,0.403870) (7,0.427612) (5,0.451599) (3,0.475769) (1,0.500000)}
\def\vSRTailTwenty{(209,0.000001) (207,0.000002) (205,0.000003) (203,0.000005) (201,0.000007) (199,0.000010) (197,0.000013) (195,0.000018) (193,0.000024) (191,0.000031) (189,0.000041) (187,0.000052) (185,0.000067) (183,0.000084) (181,0.000105) (179,0.000131) (177,0.000161) (175,0.000197) (173,0.000241) (171,0.000293) (169,0.000354) (167,0.000425) (165,0.000508) (163,0.000605) (161,0.000716) (159,0.000845) (157,0.000993) (155,0.001163) (153,0.001356) (151,0.001576) (149,0.001827) (147,0.002110) (145,0.002430) (143,0.002790) (141,0.003195) (139,0.003648) (137,0.004154) (135,0.004718) (133,0.005344) (131,0.006040) (129,0.006808) (127,0.007656) (125,0.008591) (123,0.009617) (121,0.010742) (119,0.011975) (117,0.013321) (115,0.014788) (113,0.016384) (111,0.018117) (109,0.019995) (107,0.022027) (105,0.024220) (103,0.026585) (101,0.029129) (99,0.031861) (97,0.034790) (95,0.037926) (93,0.041275) (91,0.044847) (89,0.048654) (87,0.052699) (85,0.056993) (83,0.061546) (81,0.066363) (79,0.071453) (77,0.076823) (75,0.082479) (73,0.088427) (71,0.094674) (69,0.101225) (67,0.108084) (65,0.115256) (63,0.122744) (61,0.130549) (59,0.138678) (57,0.147126) (55,0.155897) (53,0.164991) (51,0.174405) (49,0.184138) (47,0.194188) (45,0.204549) (43,0.215217) (41,0.226188) (39,0.237453) (37,0.249004) (35,0.260837) (33,0.272938) (31,0.285298) (29,0.297910) (27,0.310757) (25,0.323828) (23,0.337111) (21,0.350591) (19,0.364253) (17,0.378083) (15,0.392063) (13,0.406178) (11,0.420411) (9,0.434744) (7,0.449159) (5,0.463639) (3,0.478165) (1,0.492718)}


\begin{document}
\maketitle

\begin{abstract}
\noindent
Рассматривается непараметрическая (distribution-free) контрольная карта
Шухарта, построенная на знаково-ранговой статистике Вилкоксона
(signed-rank, SR). Документ объединяет математические основы метода,
описание программной реализации класса \texttt{SignedRankChart},
сопоставление с доступными реализациями на языке R (пакеты \texttt{npSPC},
\texttt{qcr} и самодельные скрипты) и вычислительный эксперимент,
включающий перекрёстную проверку на идентичных данных, проверку
достигнутого уровня ложных тревог и оценку вероятности обнаружения сдвига.
Отдельно разобраны расхождения между консервативным правилом границ из книги
и асимптотическим режимом пакета \texttt{npSPC}.
\end{abstract}

% ================================================================== %
\section*{Сокращения}
\addcontentsline{toc}{section}{Сокращения}
% ================================================================== %
Сокращения приведены в исходном (английском) написании, принятом в
литературе по статистическому управлению процессами; пояснения даны
по-русски.

\begin{description}[leftmargin=2.8cm,itemsep=1pt]
  \item[SR] signed-rank statistic --- знаково-ранговая статистика (сумма
        рангов модулей со знаками отклонений от цели);
  \item[SN] sign statistic --- знаковая статистика (для сравнения);
  \item[CL] center line --- центральная линия;
  \item[UCL, LCL] upper/lower control limit --- верхняя и нижняя контрольные
        границы;
  \item[FAR] false alarm rate --- вероятность ложной тревоги на одной
        подгруппе;
  \item[ARL] average run length --- средняя длина серии;
  \item[MRL] median run length --- медианная длина серии;
  \item[SD] standard deviation --- среднеквадратичное отклонение;
  \item[IC, OOC] in-control, out-of-control --- управляемое и неуправляемое
        состояние процесса;
  \item[$\mu_0$] целевое (референтное) значение центра процесса;
  \item[DP] dynamic programming --- динамическое программирование
        (свёртка распределения).
\end{description}

% ================================================================== %
\section{Математические основы реализации}
\label{sec:math}
% ================================================================== %

\subsection{Знаково-ранговая статистика}
\label{sec:sr-stat}

Знаково-ранговая карта Шухарта восходит к знаково-ранговому критерию
Вилкоксона \cite{Wilcoxon1945,Bakir2004} и переносу его на контроль
процесса \cite{Gibbons2010,Chakraborti2019}. Пусть измерение в подгруппе
объёма $n$ есть $X_1,\dots,X_n$, а $\mu_0$ --- целевое значение центра
процесса. Определим ранги модулей отклонений от цели:
\begin{equation}
R_j=\operatorname{rank}\bigl(|X_j-\mu_0|\bigr),\qquad j=1,\dots,n,
\label{eq:rank}
\end{equation}
и образуем статистику
\begin{equation}
\SR_i=\sum_{j=1}^{n}\operatorname{sign}(X_j-\mu_0)\,R_j
      =\sum_{R_j:\ X_j>\mu_0} R_j-\sum_{R_j:\ X_j<\mu_0} R_j
      =2T^+-S,
\label{eq:sr}
\end{equation}
где $T^+=\sum_{j:\,X_j>\mu_0} R_j$ --- сумма рангов положительных
отклонений, а
\begin{equation}
S=\frac{n(n+1)}{2}
\label{eq:S}
\end{equation}
--- максимально возможная сумма рангов. Статистика $\SR$ принимает
значения $\{-S,-S+2,\dots,S\}$ с шагом $2$, то есть той же чётности, что и
$S$ (при этом $S$ и $n$ связаны через $n\equiv S\!\pmod 2$ нелинейно).

\paragraph{Ключевое допущение (симметрия распределения).} В отличие от
знаковой карты, для которой достаточно непрерывности $F$, знаково-ранговая
карта требует, чтобы в управляемом состоянии распределение $F$ было
\emph{непрерывным и симметричным} относительно $\mu_0$ (класс Case~K в
терминологии книги \cite[гл.~4]{Chakraborti2019}). Именно симметрия
обеспечивает одинаковую вероятность знаков при заданном модуле отклонения,
а значит --- равновероятность всех $2^n$ знаковых наборов и
distribution-free свойство $\SR$. Для несимметричных распределений карта
неприменима; в таких случаях следует использовать знаковую карту (SN) или
другие методы.

\subsection{Распределение и моменты}
\label{sec:sr-dist}

При выполнении допущения симметрии и при $\mu_0$, равном центру, знаки
$s_j=\operatorname{sign}(X_j-\mu_0)$ независимы и равновероятны, а
$\SR=\sum_{j=1}^{n}s_j\,j$. Поэтому
\begin{equation}
P(\SR=v)=\frac{1}{2^n}\,\#\Bigl\{s\in\{-1,+1\}^n:\ \textstyle\sum_j s_j\,j=v\Bigr\}.
\label{eq:dist}
\end{equation}
Прямое перечисление всех $2^n$ знаковых наборов быстро становится
неосуществимым ($n\ge22$). В реализации распределение строится
\emph{динамическим программированием} (свёрткой): обозначив $c_j(v)$ число
наборов из первых $j$ рангов с суммой $v$, имеем
\begin{equation}
c_0(0)=1,\qquad
c_j(v)=c_{j-1}(v-j)+c_{j-1}(v+j),\qquad
P(\SR=v)=\frac{c_n(v)}{2^n}.
\label{eq:dp}
\end{equation}
Сложность $O(n\cdot S)$ вместо $O(2^n n)$; при $n=30$ ($S=465$) это доли
миллисекунды. Моменты распределения известны точно
\cite[с.~200, табл.~4.3]{Chakraborti2019}:
\begin{equation}
E[\SR]=0,\qquad
\SD(\SR)=\sqrt{\frac{n(n+1)(2n+1)}{6}}.
\label{eq:sd}
\end{equation}
Например, $\SD(\SR)\big|_{n=16}=\sqrt{1496}\approx38.68$.

\subsection{Контрольные границы и правило сигнала}
\label{sec:sr-limits}

Границы карты симметричны и задаются одной константой $a$:
\begin{equation}
\mathrm{UCL}=+a,\qquad \mathrm{CL}=0,\qquad \mathrm{LCL}=-a,
\qquad 1\le a\le S .
\label{eq:limits}
\end{equation}
Процесс признаётся неуправляемым (сигнал), если
\begin{equation}
|\SR_i|\ge a .
\label{eq:signal}
\end{equation}
При симметричном распределении $\SR$ двусторонняя вероятность ложной
тревоги равна
\begin{equation}
\FAR(a)=2\,P(\SR\ge a),\qquad
\ARL=\frac{1}{\FAR},\qquad
\MRL=\left\lceil \frac{\ln 0.5}{\ln(1-\FAR)}\right\rceil .
\label{eq:far}
\end{equation}
Отраслевой стандарт --- $\FARz=0.0027$, что соответствует трёхсигмовым
границам нормального распределения.

\subsection{Методы выбора константы $a$}
\label{sec:methods}

Книжное правило \cite[с.~198, Note~4.2]{Chakraborti2019} ---
\emph{точное} (консервативное):
\begin{equation}
a_{\mathrm{exact}}=\min\bigl\{\,a:\ 1\le a\le S,\ a\equiv S\ (\mathrm{mod}\ 2),\
P(\SR\ge a)\le \FARz/2 \,\bigr\},
\label{eq:cons}
\end{equation}
с запасным значением $a=S$, если условие не выполняется ни при каком $a$.
Такое $a$ всегда существует как значение $\SR$, а достигнутый FAR не
превосходит $\FARz$: $\FAR(a_{\mathrm{exact}})=2P(\SR\ge a_{\mathrm{exact}})
\le \FARz$.

Библиотека дополнительно реализует два вспомогательных режима, которых в
книге нет и которые служат для сравнения и обучения.

\paragraph{Нормальная аппроксимация.} Граница берётся по
нормальному приближению распределения $\SR$:
\begin{equation}
a_{\mathrm{norm}}=\min\bigl\{\,a:\ 1\le a\le S,\ a\equiv S\ (\mathrm{mod}\ 2),\
a\ge z_{1-\FARz/2}\,\SD(\SR)\,\bigr\},
\label{eq:sym}
\end{equation}
где $z_{1-\FARz/2}=z_{0.99865}\approx2.9999$ --- квантиль стандартного
нормального закона, $\SD(\SR)$ --- из \eqref{eq:sd}. Это асимптотический по
природе метод: при малых $n$ он совпадает с точным, а с ростом $n$
может давать \emph{более широкую} границу, поэтому достигнутый FAR не
обязан быть не больше $\FARz$. Режим воспроизводит границы
\texttt{npSPC::Limits\_asymptotical} (см.~разделы~\ref{sec:npspc}
и~\ref{sec:disc}).

\paragraph{Ближайший.} Минимизируется расстояние до номинального уровня
$\alpha$ (по умолчанию $\alpha=0.05$):
\begin{equation}
a_{\mathrm{near}}=\arg\min_{1\le a\le S}\bigl|\,2P(\SR\ge a)-\alpha\,\bigr|.
\label{eq:near}
\end{equation}
Гарантии $\FAR\le\FARz$ нет; режим приводится исключительно для полноты
картины и в производстве не применяется.

\begin{table}[htbp]\centering
\caption{Режимы выбора константы $a$.}
\label{tab:srmodes}
\small
\begin{tabular}{@{}p{2.5cm} p{3.1cm} p{7.4cm} p{1.5cm}@{}}
\toprule
Режим & Целевой риск & Правило & Чёт\-ность \\
\midrule
\texttt{exact} & $\FARz/2=0.00135$ & наименьшее достижимое $a$ с
$P(\SR\ge a)\le\FARz/2$ & да\\
\texttt{normal} & $z_{1-\FARz/2}\SD(\SR)$ & наименьшее достижимое
$a\ge z_{1-\FARz/2}\SD(\SR)$ & да\\
\texttt{nearest} & $\alpha=0.05$ & $\arg\min_a|2P(\SR\ge a)-\alpha|$ & нет\\
\bottomrule
\end{tabular}
\normalsize
\end{table}

\begin{figure}[htbp]\centering
\begin{tikzpicture}
\begin{axis}[
  width=0.86\textwidth, height=6.4cm,
  xlabel={$n$}, ylabel={константа $a$},
  xmin=2, xmax=20, ymin=0, ymax=170,
  legend pos=north west, legend cell align=left,
  title={Границы по режимам}]
\addplot+[mark=*,mark size=1.3pt,color=accent] coordinates {\vSRACons};
\addlegendentry{\texttt{exact} (книга)}
\addplot+[mark=square*,mark size=1.3pt,color=third] coordinates {\vSRASym};
\addlegendentry{\texttt{normal} (норм.\ аппр.)}
\addplot+[mark=triangle,mark size=1.5pt,color=alert] coordinates {\vSRANear};
\addlegendentry{\texttt{nearest} ($\alpha=0.05$)}
\end{axis}
\end{tikzpicture}
\caption{Зависимость константы $a$ от объёма подгруппы $n$.
\texttt{exact} и \texttt{normal} совпадают при $n\le10$ и
расходятся далее; \texttt{nearest} систематически даёт самые узкие границы.}
\label{fig:modesa}
\end{figure}

\begin{figure}[htbp]\centering
\begin{tikzpicture}
\begin{axis}[
  width=0.86\textwidth, height=6.2cm,
  xlabel={$n$}, ylabel={достигнутый $\FAR$},
  xmin=2, xmax=20, ymin=0, ymax=0.06,
  legend pos=north east, legend cell align=left,
  title={Достигнутый уровень ложных тревог}]
\addplot+[mark=*,mark size=1.3pt,color=accent] coordinates {\vSRFARCons};
\addlegendentry{\texttt{exact}}
\addplot+[mark=square*,mark size=1.3pt,color=third] coordinates {\vSRFARSym};
\addlegendentry{\texttt{normal}}
\addplot+[mark=triangle,mark size=1.5pt,color=alert] coordinates {\vSRFARNear};
\addlegendentry{\texttt{nearest}}
\addplot[dashed,color=softgray,domain=2:20] {0.0027};
\addlegendentry{$\FARz=0.0027$}
\end{axis}
\end{tikzpicture}
\caption{Достигнутый $\FAR$ по режимам. Точный режим держится ниже
$\FARz$; нормальная аппроксимация при $n\ge11$ уходит ниже (граница шире);
\texttt{nearest} колеблется около $\alpha=0.05$.}
\label{fig:modesfar}
\end{figure}

\begin{table}[htbp]\centering
\caption{Сравнение режимов, $n=2..20$ ($\FARz=0.0027$, $\alpha=0.05$).}
\label{tab:srmodecmp}
\small
\begin{tabular}{@{}rrrrrrr@{}}
\toprule
$n$ & $a_{\mathrm{exact}}$ & $\FAR_{\mathrm{exact}}$ & $a_{\mathrm{norm}}$
& $\FAR_{\mathrm{norm}}$ & $a_{\mathrm{near}}$ & $\FAR_{\mathrm{near}}$\\
\midrule
\vSRModeTable
\bottomrule
\end{tabular}
\normalsize
\end{table}

\paragraph{Пример ($n=14$, $S=105$).} Рассмотрим три режима подробнее.
\begin{itemize}
\item \emph{Точный.} Условие $P(\SR\ge a)\le0.00135$ впервые
выполняется при $a=91$: $P(\SR\ge91)=0.001160$, тогда как $a=93$ даёт
$0.000854$, а $a=89$ уже $0.001526>0.00135$. Наименьшее достижимое ---
$a=91$ ($\FAR=0.002319$, $\ARL\approx431$).
\item \emph{Нормальная аппроксимация.} $\SD(\SR)=\sqrt{1015}\approx31.86$, откуда
$z_{1-0.00135}\cdot31.86\approx95.6$; ближайшее достижимое значение (той же
чётности, что и $S=105$) --- $a=97$ ($\FAR=0.000854$, $\ARL\approx1171$).
Асимптотическая граница здесь шире точной.
\item \emph{Ближайший.} Минимум $|2P(\SR\ge a)-0.05|$ достигается при
$a=62$ (и эквивалентном $a=63$): $2P\approx0.0494$.
\end{itemize}

\begin{figure}[htbp]\centering
\begin{tikzpicture}
\begin{axis}[
  width=0.86\textwidth, height=6.0cm,
  xlabel={$a$}, ylabel={$2\,P(|\SR|\ge a)$},
  xmin=40, xmax=80, ymin=0, ymax=0.24,
  title={Критерий \texttt{nearest}, $n=14$}]
\addplot+[mark=*,mark size=1.2pt,color=accent] coordinates {\vSRNearAbs14};
\addplot+[mark=square*,mark size=2.6pt,color=alert] coordinates {\vSRNearMin14};
\addplot[dashed,color=softgray,domain=40:80] {0.05};
\end{axis}
\end{tikzpicture}
\caption{Функция $2P(|\SR|\ge a)$ для $n=14$. Ромбом отмечен минимум
отклонения от $\alpha=0.05$ (значения $a=62$ и $a=63$ дают одинаковую
вероятность, так как $\SR$ чётности $S$).}
\label{fig:modesnear}
\end{figure}

Для иллюстрации распределения на рисунке~\ref{fig:srpmf} приведена
дискретная плотность $\SR$ при $n=5$, а на рисунке~\ref{fig:srtail} ---
односторонние хвосты $P(\SR\ge a)$ при различных $n$.

\begin{figure}[htbp]\centering
\begin{tikzpicture}
\begin{axis}[
  width=0.8\textwidth, height=5.8cm,
  xlabel={значение $\SR$}, ylabel={$P(\SR=v)$},
  xmin=-16, xmax=16, ymin=0, ymax=0.11,
  ybar, bar width=4pt,
  title={Распределение $\SR$ в IC, $n=5$ ($2^5=32$ набора)}]
\addplot+[mark=none,color=accent,fill=accent!45] coordinates {\vSRpmfFive};
\end{axis}
\end{tikzpicture}
\caption{Распределение знаково-ранговой статистики при $n=5$: симметрично,
сосредоточено на значениях чётности $S=15$.}
\label{fig:srpmf}
\end{figure}

\begin{figure}[htbp]\centering
\begin{tikzpicture}
\begin{axis}[
  width=0.86\textwidth, height=6.2cm,
  xlabel={$a$}, ylabel={$P(\SR\ge a)$},
  xmode=log, log basis x=10,
  ymin=0, ymax=0.5, legend pos=north east, legend cell align=left]
\addplot+[mark=*,mark size=1pt,color=accent] coordinates {\vSRTailFive};
\addlegendentry{$n=5$}
\addplot+[mark=square*,mark size=1pt,color=third] coordinates {\vSRTailTen};
\addlegendentry{$n=10$}
\addplot+[mark=triangle,mark size=1pt,color=alert] coordinates {\vSRTailFourteen};
\addlegendentry{$n=14$}
\addplot+[mark=diamond*,mark size=1pt,color=softgray] coordinates {\vSRTailTwenty};
\addlegendentry{$n=20$}
\end{axis}
\end{tikzpicture}
\caption{Односторонний хвост $P(\SR\ge a)$ (логарифмическая шкала $a$).
С ростом $n$ хвост становится круче, а достижимая граница при
фиксированном $\FARz$ --- всё ближе по относительной вероятности к цели.}
\label{fig:srtail}
\end{figure}

\paragraph{Листинг выбора $a$.} Ниже приведён фрагмент
\texttt{src/nonparametric\_charts/charts/signed\_rank\_chart.py}.

\begin{lstlisting}[language=Python]
@classmethod
def _calc_a_exact(cls, n, far0, dist_func):
    target = far0 / 2
    S = n * (n + 1) // 2
    start = 1 if S % 2 == 1 else 2
    for a in range(start, S + 1, 2):
        if dist_func(n, a) <= target:
            return a, 2.0 * dist_func(n, a)
    return S, 2.0 * dist_func(n, S)   # fallback: widest limits

@classmethod
def _calc_a_normal(cls, n, far0, dist_func):
    S = n * (n + 1) // 2
    z = stats.norm.ppf(1.0 - far0 / 2.0)
    sd = (n * (n + 1) * (2 * n + 1) / 6.0) ** 0.5
    a0 = int(np.ceil(z * sd))
    a = a0 if a0 % 2 == S % 2 else a0 + 1
    a = min(max(a, 1), S)
    return a, 2 * dist_func(n, a)
\end{lstlisting}

\subsection{Контрольные границы и дискретность}
\label{sec:sr-discrete}

Из--за дискретности $\SR$ достичь номинального $\FARz$ точно, как правило,
не удаётся: приходится либо округлять границу вверх (расширяя её и занижая
достигнутый FAR), либо вниз (завышая FAR). Книжное правило
\eqref{eq:cons} выбирает наименьшее достижимое $a$, при котором
односторонний хвост ещё не превышает $\FARz/2$, --- это округление вверх и
гарантия $\FAR\le\FARz$. Нормальная аппроксимация \eqref{eq:sym} такой
гарантии не даёт, но при больших $n$ становится точнее по вероятности.

% ================================================================== %
\section{Программная реализация}
\label{sec:impl}
% ================================================================== %

\subsection{Общая структура и зависимости}
\label{sec:impl-struct}

Знаково-ранговая карта реализована классом \texttt{SignedRankChart}
(\path{src/nonparametric_charts/charts/signed_rank_chart.py}), наследующим базовый класс
\texttt{BaseControlChart} (\path{src/nonparametric_charts/base/base_chart.py}). Дополнительно
используются модули \path{src/nonparametric_charts/utils/dataloader.py},
\path{src/nonparametric_charts/utils/preprocessor.py}, \path{src/nonparametric_charts/utils/plotting.py} и единый
стиль графиков \path{src/nonparametric_charts/utils/plot_style.py}. Импорты абсолютные
(\texttt{from nonparametric\_charts.charts...}), запуск --- из корня репозитория; требования
--- Python $\ge3.14$ и библиотеки \texttt{pandas}, \texttt{numpy},
\texttt{scipy}, \texttt{matplotlib}. Коэффициенты
\path{coefficients.json} относятся к параметрическим картам и здесь не
используются.

Входные данные --- таблица \texttt{pandas.DataFrame}: строка есть подгруппа,
столбец --- измерение, параметр \texttt{sample\_size} задаёт число
столбцов (объём подгруппы $n$). Целевое значение $\mu_0$ задаётся явно или
оценивается медианой всех наблюдений.

\subsection{Класс \texttt{SignedRankChart}: статистика по подгруппам}
\label{sec:impl-stat}

Метод \texttt{\_calculate\_sr\_for\_row} для каждой подгруппы выполняет
шаги \eqref{eq:rank}--\eqref{eq:sr}: отбрасывает пропуски, вычисляет
отклонения $d_j=x_j-\mu_0$, отбрасывает (теоретически невозможные) нули,
ранжирует $|d_j|$ методом \texttt{average} (учёт ничьих при округлении),
и возвращает $\SR$, а также суммы рангов $T^+$ (положительных) и
отрицательных. Результаты собираются в таблицу \texttt{self.subgroups} с
колонками \texttt{n\_used}, \texttt{positive}, \texttt{negative},
\texttt{zero}, \texttt{sr}, \texttt{sum\_pos\_ranks},
\texttt{sum\_neg\_ranks}, \texttt{T\_plus}.

\subsection{Класс \texttt{SignedRankChart}: распределение и константа}
\label{sec:impl-dist}

Распределение $\SR$ в IC (\texttt{\_sr\_distribution}) строится
динамическим программированием по рекурсии \eqref{eq:dp} и кэшируется по
$n$. Функции \texttt{\_p\_sr\_ge\_a} и \texttt{\_p\_abs\_sr\_ge\_a}
возвращают односторонний хвост $P(\SR\ge a)$ и двустороннюю вероятность
$P(|\SR|\ge a)$. Метод \texttt{\_calculate\_a\_value} выбирает $a$
согласно режиму: \texttt{exact}, \texttt{normal} или
\texttt{nearest} (см.~листинг в~разделе~\ref{sec:methods}). Устаревшие
идентификаторы \texttt{approach='conservative'} и
\texttt{approach='symmetric'} принимаются для обратной совместимости с
предупреждением об устаревании и отображаются на \texttt{exact} и
\texttt{normal} соответственно (см.~ADR-012 в памяти проекта).

\subsection{Границы, нарушения и визуализация}
\label{sec:impl-plot}

Метод \texttt{calculate\_limits} формирует статистики по всем подгруппам,
определяет $a$ (или принимает заданное), записывает границы
$\mathrm{UCL}=+a$, $\mathrm{CL}=0$, $\mathrm{LCL}=-a$ вместе с
достигнутым $\FAR$, $\ARL$ и $\MRL$, а затем вызывает
\texttt{\_find\_violations}, реализующий правило сигнала
$|\SR_i|\ge a$ \cite[\S4.2.1.2.1, ур.~4.6, с.~198]{Chakraborti2019}.
Метод \texttt{plot} рисует карту с линией статистики, границами, линией
$\mathrm{CL}$ и (при \texttt{show\_zones=True}) выделение зон; метод
\texttt{plot\_comparison} строит семейство кривых $a$-по-$n$ для нескольких
режимов. Подписи используют русские наименования режимов из
\texttt{APPROACH\_LABELS}, причём для \texttt{normal} ---
<<Нормальная аппроксимация>>.

\subsection{Фазы I и II и вычислительная сложность}
\label{sec:impl-complexity}

Реализация поддерживает обе фазы: в фазе~I целевое значение и границы
оцениваются по данным, в фазе~II используются зафиксированные значения; для
этого \texttt{calculate\_limits} принимает явное \texttt{a\_value}, а
$\mu_0$ --- параметр конструктора. Основные затраты --- ранжирование по
подгруппам ($O(m\,n\log n)$ при $m$ подгруппах) и построение
распределения $\SR$ ($O(n\cdot S)$, однократно на каждое $n$). Благодаря
ДП-свёртке распределение доступно и при $n\ge30$, где перебор $2^n$
неосуществим.

\begin{table}[htbp]\centering
\caption{Ориентировочное время работы ($n=16$, если не указано иное).}
\label{tab:timing}
\small
\begin{tabular}{@{}l r@{}}
\toprule
Операция & Время \\
\midrule
Наш расчёт границ, 1000 раз ($n=16$) & $\approx0.39$~с \\
Правило \texttt{npSPC} exact, 1000 раз ($n=16$) & $\approx0.32$~с \\
Распределение $\SR$ ДП-свёрткой, $n=30$ & $\approx0.002$~с \\
Обработка $10^6$ подгрупп векторно & $\approx0.52$~с \\
\bottomrule
\end{tabular}
\normalsize
\end{table}

% ================================================================== %
\section{Реализации на языке R}
\label{sec:r}
% ================================================================== %

\subsection{Пакет \texttt{npSPC}}
\label{sec:npspc}

Пакет \texttt{npSPC} \cite{Nascimento2026} содержит функцию
\texttt{shewhart\_sr\,(X, mu0, far, group\_by\_col, exact, side)}.
Статистика вычисляется так же, как в \eqref{eq:sr}:
\begin{lstlisting}[language=R]
Ranks = apply(abs(X - mu0), 2, rank)
Tplus = colSums(apply((X - mu0) > 0, 2, ifelse, 1, 0) * Ranks, na.rm = T)
SR    = 2*Tplus - n*(n+1)/2
\end{lstlisting}
Существенные особенности:
\begin{itemize}
\item матрица данных ориентирована <<строка --- наблюдение внутри подгруппы,
столбец --- подгруппа>> (при \texttt{group\_by\_col=TRUE} выполняется
транспонирование);
\item параметр \texttt{far} при \texttt{side="two.sided"} делится пополам:
\texttt{alpha=far/2};
\item при \texttt{exact=FALSE} (по умолчанию) границы берутся из
\texttt{Limits\_asymptotical}:
$\mathrm{UCL}=z_{1-\alpha}\,\SD(\SR)$ --- \emph{дробное} число, без
приведения к достижимому значению $\SR$; при $n<20$ печатается
предупреждение о желательности точного режима;
\item при \texttt{exact=TRUE} вызывается \texttt{Exact\_distribution}:
строится распределение $\SR$ перебором $2^n$ наборов, затем
$\mathrm{UCL}=\min\{u:\ \mathrm{CDF}(u)\ge 1-\alpha,\ u\le S\}$ ---
это \emph{квантиль} уровня $1-\alpha$ (<<прохождение>> ступеньки CDF), а
не наименьшее $a$ с $P(\SR\ge a)\le\alpha$.
\end{itemize}
Функция возвращает строку с номинальным $\ARL_0=1/\far$, а вычисленные
границы и статистику лишь передаёт во внутренний \texttt{plot\_chart}.

\subsection{Пакет \texttt{qcr}}
\label{sec:qcr}

Пакет \texttt{qcr} ориентирован на контрольные карты, в том числе
непараметрические, однако его непараметрические средства многомерны и
основаны на глубине данных: функции \texttt{npqcs.*} работают с опорной
выборкой и глубинами \cite{Liu1995}, а функция \texttt{fdqcs.rank}
относится к функциональным/ранговым данным. Классической одномерной
знаково-ранговой карты (по статистике \eqref{eq:sr}) в пакете нет;
ближайший одномерный инструмент --- атрибутная карта числа несоответствий
\texttt{qcs.np}, не связанная с Wilcoxon-статистикой. Таким образом,
\texttt{qcr} не является конкурентом рассматриваемой реализации и
интересен лишь как источник родственных методов.

\subsection{Реализации средствами базового R}
\label{sec:baser}

Пользовательские скрипты строятся из \texttt{rank} и \texttt{colSums} и,
как правило, задают границы по нормальной аппроксимации
$\pm z_{1-\far/2}\SD(\SR)$:
\begin{lstlisting}[language=R]
SR <- 2 * colSums((X - med > 0) * apply(abs(X - med), 2, rank)) - n*(n+1)/2
a  <- ceiling(qnorm(1 - far / 2) * sqrt(n * (n + 1) * (2 * n + 1) / 6))
signal <- which(abs(SR) >= a)
\end{lstlisting}
Расчёт \path{experiments/sr/sr_custom.R} показывает, что при $n=16$
нормальная граница ($a=117$) на семь единиц шире точной книжной ($a=110$),
а при $n=40$ разрыв достигает девяти единиц; при этом из--за чётности
достигнутый FAR нормальной границы оказывается \emph{ниже} номинального
(около $0.0010$). Это качественно отличается от знаковой карты, где
нормальная аппроксимация занижает границу.

% ================================================================== %
\section{Сопоставление реализаций}
\label{sec:compare}
% ================================================================== %

Сводное сопоставление приведено в таблице~\ref{tab:compare}. Ключевые
наблюдения:
\begin{itemize}
\item наша реализация и \texttt{npSPC} дают \emph{идентичную} статистику
$\SR$ на одних и тех же данных (перекрёстная проверка,
раздел~\ref{sec:cross});
\item книжное правило \eqref{eq:cons} и точный режим \texttt{npSPC}
различаются начинкой: у нас --- точное <<наименьшее достижимое $a$
с $P(\SR\ge a)\le\far/2$>>, у пакета --- квантиль уровня $1-\far/2$; их
границы начинают расходиться уже при $n\ge10$ (раздел~\ref{sec:disc});
\item режим \texttt{normal} (нормальная аппроксимация) воспроизводит
границы \texttt{npSPC} по умолчанию (\texttt{exact=FALSE}) с точностью до
приведения к достижимому значению $\SR$;
\item \texttt{qcr} классической знаково-ранговой карты не предоставляет.
\end{itemize}

\begin{table}[htbp]\centering
\caption{Сопоставление реализаций знаково-ранговой карты.}
\label{tab:compare}
\footnotesize
\begin{tabular}{@{}p{2.7cm} p{3.4cm} p{3.7cm} p{4.2cm}@{}}
\toprule
Реализация & Статистика & Границы & Особенности \\
\midrule
\texttt{SignedRankChart} (наш) & $\SR=2T^+-S$, ранги \texttt{average} &
три режима: \texttt{exact}, \texttt{normal}, \texttt{near.} &
учёт чётности $S$; ДП-свёртка; $\mu_0$, фазы I/II \\
\texttt{npSPC} (\texttt{shewhart\_sr}) & $\SR=2T^+-S$, \texttt{rank} &
\texttt{exact} (квантиль), асимптотические (дробные) &
\texttt{far/2}; перебор $2^n$; только фаза~II \\
Базовый R (самодельный) & $\SR=2T^+-S$ & нормальная аппроксимация &
без точного распределения, без чётности \\
\texttt{qcr} & --- & --- & классической SR нет \\
\bottomrule
\end{tabular}
\normalsize
\end{table}

% ================================================================== %
\section{Вычислительный эксперимент}
\label{sec:exp}
% ================================================================== %

Все расчёты воспроизводимы: точные вероятности считаются ДП-свёрткой,
проверочные опыты --- методом Монте-Карло с фиксированным состоянием
генератора. Скрипты и инструкции --- в \path{experiments/sr/};
среда описана в \path{experiments/00_ENVIRONMENT.txt}.

\subsection{Перекрёстная проверка на идентичных данных}
\label{sec:cross}

По методике \path{sr_cross_validation.R} сгенерирована матрица
$n=16$, $m=3000$ подгрупп ($N(0,1)$); начиная с подгруппы~2001 введён
сдвиг центра на $1$ SD. Данные переданы в Python
(\path{sr_cross_check.py}). Результаты:
\begin{itemize}
\item знаково-ранговая статистика $\SR$ полностью совпадает у нашей
реализации и у \texttt{npSPC} (\texttt{True}) --- ранжирование и формула
\eqref{eq:sr} идентичны;
\item точная граница \texttt{npSPC} при $\far=0.0027$ равна $a=108$, наша
(режим \texttt{exact}) --- $a=110$;
\item на управляющем участке (2000 подгрупп) граница $108$ дала 5 ложных
тревог, граница $110$ --- 3, что согласуется с их теоретическими $\FAR$
($0.0034$ и $0.0027$);
\item ничьих (\texttt{X==0}) нет --- данные непрерывны.
\end{itemize}

\subsection{Проверка уровня ложных тревог}
\label{sec:far}

Опыт Монте-Карло ($n=16$, $m=5\cdot10^5$ управляемых подгрупп)
подтверждает теорию (таблица~\ref{tab:far}). Эмпирические частоты
согласуются с точными в пределах стандартной ошибки
($\sqrt{\FARz(1-\FARz)/m}\approx7\cdot10^{-5}$).

\begin{table}[htbp]\centering
\caption{Уровень ложных тревог, $n=16$, $m=5\cdot10^5$.}
\label{tab:far}
\small
\begin{tabular}{@{}l r r r r@{}}
\toprule
Режим & $a$ & $\FAR$ (точный) & $\FAR$ (опыт) & $\ARL$ (опыт) \\
\midrule
точный (книга) & 110 & 0.002686 & 0.002778 & 360.0 \\
нормальная аппроксимация & 118 & 0.001007 & 0.001014 & 986.2 \\
\texttt{npSPC} exact & 108 & 0.003357 & 0.003444 & 290.4 \\
\bottomrule
\end{tabular}
\normalsize
\end{table}

\subsection{Мощность при сдвиге}
\label{sec:power}

Вероятность обнаружения сдвига центра на $\delta\cdot\SD(\SR)$ вычисляется
двумя способами: напрямую методом Монте-Карло (моделирование данных) и
свёрткой $\SR=\sum_j s_j j$, где $s_j$ --- знак $j$-й по модулю порядковой
статистики с вероятностью $q_j$ (формула через плотность нормального
распределения). Свёртка даёт точное распределение при гипотезе
независимости ранговых знаков; это приближение, но оно согласуется с
Монте-Карло в пределах 1--2\,\% (таблица~\ref{tab:power}). Границы
$a=110$ (точная) и $a=108$ (\texttt{npSPC}) различаются по мощности
мало: при $\delta=1$ SD выигрыш узкой границы --- около 4 процентных
пунктов, но ценой повышенного FAR.

\begin{table}[htbp]\centering
\caption{Вероятность сигнала и $\ARL$, $n=16$ (свёртка).}
\label{tab:power}
\small
\begin{tabular}{@{}r r r r r@{}}
\toprule
$\delta$ & $\Phi(\delta)$ & $P$ ($a=110$) & $\ARL$ ($a=110$)
& $P$ ($a=108$) \\
\midrule
0.0 & 0.500 & 0.002686 & 372.3 & 0.003357 \\
0.5 & 0.691 & 0.086417 & 11.57 & 0.100213 \\
1.0 & 0.841 & 0.623783 & 1.60 & 0.662585 \\
1.5 & 0.933 & 0.980478 & 1.02 & 0.985944 \\
2.0 & 0.977 & 0.999897 & 1.00 & 0.999948 \\
\bottomrule
\end{tabular}
\normalsize
\end{table}

\begin{figure}[htbp]\centering
\begin{tikzpicture}
\begin{axis}[
  width=0.86\textwidth, height=6.2cm,
  xlabel={сдвиг $\delta$ (в $\SD$)}, ylabel={$P(\text{сигнал})$},
  xmin=0, xmax=2, ymin=0, ymax=1.05,
  legend pos=south east, legend cell align=left]
\addplot+[mark=*,mark size=1.4pt,color=accent] coordinates {\vSROCBk};
\addlegendentry{книга, $a=110$}
\addplot+[mark=square*,mark size=1.4pt,color=alert] coordinates {\vSROCNp};
\addlegendentry{\texttt{npSPC} exact, $a=108$}
\end{axis}
\end{tikzpicture}
\caption{Кривые обнаружения сдвига при $n=16$. Более узкая граница пакета
быстрее реагирует, но при повышенном FAR в управляемом состоянии.}
\label{fig:oc}
\end{figure}

\subsection{Анализ расхождений}
\label{sec:disc}

Причина расхождения книжного правила и точного режима \texttt{npSPC} ---
разная трактовка дискретности. Книжное правило \eqref{eq:cons} берёт
\emph{наименьшее достижимое} $a$, для которого односторонний хвост ещё не
превышает $\FARz/2$ (округление вверх, гарантия $\FAR\le\FARz$). Пакет
\texttt{npSPC} при \texttt{exact=TRUE} берёт \emph{квантиль} уровня
$1-\FARz/2$ --- точку, где ступенчатая CDF впервые пересекает уровень
(округление, эквивалентное выбору значения, при котором хвост может
\emph{превышать} $\FARz/2$). Поэтому границы пакета уже, а достигнутый FAR
может быть выше номинального (таблица~\ref{tab:limits}).

\begin{table}[htbp]\centering
\caption{Границы и достигнутый FAR по правилам.}
\label{tab:limits}
\small
\begin{tabular}{@{}r r r r r r r@{}}
\toprule
$n$ & $S$ & книга $a$ & FAR & \texttt{npSPC} exact $a$ & его FAR
& $\lceil z\SD\rceil$ \\
\midrule
10 & 55 & 55 & 0.001953 & 53 & 0.003906 & 59$^{\dagger}$ \\
12 & 78 & 72 & 0.002441 & 70 & 0.003418 & 77 \\
14 & 105 & 91 & 0.002319 & 89 & 0.003052 & 96 \\
16 & 136 & 110 & 0.002686 & 108 & 0.003357 & 117 \\
18 & 171 & 133 & 0.002335 & 131 & 0.002808 & 138 \\
20 & 210 & 156 & 0.002325 & 154 & 0.002712 & 161 \\
\bottomrule
\end{tabular}
\normalsize
\\[2pt]
{\footnotesize $^{\dagger}$ Последний столбец --- \emph{дробная}
асимптотическая граница $z\cdot\SD(\SR)$, округлённая вверх. Наша
реализация затем приводит её к чётности $S$ и зажимает $a\le S$, получая
эффективную границу \texttt{normal}: $55$, $78$, $97$, $118$, $139$,
$162$. Пакет \texttt{npSPC} (режим \texttt{exact=FALSE}) такого зажима не
делает, поэтому при $n=10$ его граница $59>S=55$ недостижима и карта не
подаёт сигналов вовсе.}
\end{table}

На рисунке~\ref{fig:limits} сопоставлены границы книжного правила, точного
режима \texttt{npSPC}, нормальной аппроксимации и нашей реализации. При
$n\le9$ книжное правило, точный режим \texttt{npSPC} и \texttt{normal}
совпадают (дискретность вынуждает $a=S$, то есть карта не сигналит при
номинале); расхождения возникают при $n\ge10$. Асимптотическая граница
пакета уже при малых $n$ превосходит $S$ и потому недостижима.

\begin{figure}[htbp]\centering
\begin{tikzpicture}
\begin{axis}[
  width=0.88\textwidth, height=6.4cm,
  xlabel={$n$}, ylabel={константа $a$},
  xmin=2, xmax=20, ymin=0, ymax=220,
  legend pos=north west, legend cell align=left]
\addplot+[mark=*,mark size=1.2pt,color=accent] coordinates {\vSRACons};
\addlegendentry{книга (\texttt{exact})}
\addplot+[mark=square*,mark size=1.2pt,color=third] coordinates {\vSRASym};
\addlegendentry{\texttt{normal} (норм.\ аппр.)}
\addplot+[mark=triangle,mark size=1.4pt,color=alert] coordinates {\vSRANpExact};
\addlegendentry{\texttt{npSPC} exact}
\addplot+[mark=diamond*,mark size=1.2pt,color=softgray] coordinates {\vSRANpAsym};
\addlegendentry{\texttt{npSPC} asymptotical}
\end{axis}
\end{tikzpicture}
\caption{Границы $a$ по правилам. Книжное правило лежит между точным
режимом \texttt{npSPC} (снизу) и нормальной аппроксимацией (сверху);
асимптотическая граница пакета при малых $n$ уходит выше $S$ и становится
недостижимой.}
\label{fig:limits}
\end{figure}

\paragraph{Замечание об эффективной границе.} Асимптотическая граница
\texttt{npSPC} --- дробное число, поэтому при сравнении с целочисленной
статистикой она действует как ближайшее целое сверху. Так, при $n=16$
$z\cdot\SD\approx116.03$ эквивалентно порогу $117$, а поскольку значения
$\SR$ имеют чётность $S=136$, условие $|\SR|\ge117$ равносильно
$|\SR|\ge118$. Именно поэтому режим \texttt{normal} и
\texttt{npSPC} (\texttt{exact=FALSE}) дают одинаковое множество сигналов,
хотя формально границы ($118$ и $116.03$) различны.

\subsection{Сводный эксперимент по методам}
\label{sec:modes}

На одних и тех же данных ($m=3000$, сдвиг $1$ SD с подгруппы~2001) для
$n=14,16,19$ сопоставлены все режимы (таблица~\ref{tab:modes-exp}). Число
ложных тревог на управляющем участке и задержка обнаружения сдвига
подтверждают теоретические $\FAR$.

\begin{table}[htbp]\centering
\caption{Метрики режимов, $n=16$ ($m=3000$, сдвиг с подгруппы~2001).}
\label{tab:modes-exp}
\small
\begin{tabular}{@{}l r r r r@{}}
\toprule
Режим & граница & $\FAR$ теор. & ложн.\ тревог & задержка \\
\midrule
наша: \texttt{exact} & 110 & 0.002686 & 5 & 1 \\
наша: \texttt{normal} & 118 & 0.001007 & 3 & 4 \\
наша: \texttt{nearest} & 75 & 0.050659 & 98 & 1 \\
нормальная аппроксимация & 116.03 & 0.001007 & 3 & 4 \\
\texttt{npSPC} exact, \texttt{far}=0.0027 & 108 & 0.003357 & 8 & 1 \\
\texttt{npSPC} exact, \texttt{far}=0.00135 & 114 & 0.001678 & 4 & 4 \\
\texttt{npSPC} асимпт., \texttt{far}=0.0027 & 116.03 & 0.001007 & 3 & 4 \\
\bottomrule
\end{tabular}
\normalsize
\end{table}

\begin{table}[htbp]\centering
\caption{Границы по режимам для $n=14,16,19$.}
\label{tab:modes-alln}
\small
\begin{tabular}{@{}l r r r@{}}
\toprule
Режим & $n=14$ & $n=16$ & $n=19$ \\
\midrule
наша \texttt{exact} & 91 & 110 & 144 \\
наша \texttt{normal} & 97 & 118 & 150 \\
наша \texttt{nearest} & 62 & 75 & 97 \\
\texttt{npSPC} exact (0.0027) & 89 & 108 & 142 \\
\texttt{npSPC} асимпт.\ (0.0027) & 95.58 & 116.03 & 149.10 \\
\texttt{npSPC} exact (0.00135) & 93 & 114 & 150 \\
\bottomrule
\end{tabular}
\normalsize
\end{table}

\paragraph{Что лучше.} Для производственного применения естественен
точный режим \eqref{eq:cons}: он единственный гарантирует
$\FAR\le\FARz$ и воспроизводит книжное правило. Режим \texttt{nearest} с
$\alpha=0.05$ непригоден (в СУП это соответствует $\ARL\approx20$).
Нормальная аппроксимация (\texttt{normal}) имеет смысл при больших $n$
как быстрый расчёт и как средство воспроизвести поведение
\texttt{npSPC} по умолчанию. Точный режим \texttt{npSPC} занимает
промежуточное положение: он чуть агрессивнее книжного (границы уже, FAR
может превышать номинал).

% ================================================================== %
\section{Результаты и рекомендации}
\label{sec:results}
% ================================================================== %

\subsection{Преимущества рассматриваемой реализации}
\label{sec:pros}

Реализация \texttt{SignedRankChart} даёт точную, distribution-free
статистику \eqref{eq:sr} и совпадает с эталонным пакетом \texttt{npSPC} по
значениям $\SR$ на идентичных данных, что подтверждено перекрёстной
проверкой. Распределение статистики строится динамическим
программированием \eqref{eq:dp}, поэтому доступно при $n\ge30$, где
перебор $2^n$ в пакете становится неосуществимым; при этом результаты
эквивалентны полному перечислению. Книжное точное правило
гарантирует непревышение номинального уровня ложных тревог и снабжено
запасным значением $a=S$ для малых $n$, где условие $\FAR\le\FARz$
недостижимо. Реализованы обе фазы, оценка целевого значения по данным и
явное задание $\mu_0$, а также визуализация с зонами и сравнительные
графики по $n$. Точный, асимптотический и <<ближайший>> режимы
позволяют изучать компромисс между ложными тревогами и чувствительностью;
единый стиль графиков и поддержка устаревших идентификаторов
обеспечивают совместимость и воспроизводимость. Разумно устроен учёт
пропусков и ничьих: нулевые и отсутствующие отклонения не искажают
ранжирование.

\subsection{Направления развития}
\label{sec:future}

Основное ограничение карты --- требование симметрии распределения; для
несимметричных процессов необходима знаковая карта или иные подходы, и
полезно добавить автоматическую диагностику симметрии с рекомендацией
выбора метода. Целесообразно расширить набор режимов: добавить
<<квантильный>> вариант с настраиваемым уровнем в духе
\texttt{npSPC::Exact\_distribution}, чтобы сравнивать точное и
квантильное округление, а также поддержать точный режим для больших $n$
через нормальную аппроксимацию с контролем погрешности. Перспективны
расширения на другие знаково-ранговые статистики (знаково-ранговые суммы
для сдвига масштаба и знаково-ранговые карты с оцениванием параметров по
фазе~I), на схемы с переменным объёмом подгруппы и на карты
EWMA/CUSUM, построенные на $\SR$. Для оценки точности границ при конечных
$n$ полезна таблица достигнутых $\FAR$/$\ARL$ в приложении. Наконец,
автоматизация вычислительного эксперимента (единый реестр сценариев с
повторным прогоном при изменениях) снизила бы трудозатраты и риск
рассинхронизации отчёта и кода.

% ================================================================== %
\section{Заключение}
\label{sec:conclusion}
% ================================================================== %

Знаково-ранговая карта Шухарта --- непараметрический (distribution-free)
инструмент контроля сдвига центра, требующий симметрии распределения
процесса. Построена и верифицирована реализация \texttt{SignedRankChart};
показано, что её статистика совпадает с эталонной реализацией
\texttt{npSPC}, а распределение эффективно вычисляется динамическим
программированием. Установлено, что книжное точное правило и
квантильный точный режим \texttt{npSPC} различаются трактовкой
дискретности: первое гарантирует $\FAR\le\FARz$, второй может допускать
$\FAR>\FARz$. Режим \texttt{normal} (нормальная аппроксимация)
воспроизводит поведение \texttt{npSPC} по умолчанию и совпадает с ним по
эффективной границе. Для производственного применения рекомендуется
точное правило; остальные режимы служат целям сравнения и
обучения. Направления развития связаны с расширением семейства статистик и
схем управления, диагностикой симметрии и автоматизацией экспериментов.

% ================================================================== %
\appendix
\section{Приложения}
\label{sec:appendix}
% ================================================================== %

\paragraph{Воспроизведение.} Все результаты получены скриптами из
\path{experiments/sr/}; порядок запуска описан в файлах
\texttt{01}--\texttt{06} той же папки, среда --- в
\path{experiments/00_ENVIRONMENT.txt}. Перекрёстная проверка:
\begin{lstlisting}[language=bash]
R_LIBS=$RENV/lib/R/library $RENV/bin/Rscript experiments/sr/sr_cross_validation.R
poetry run python experiments/sr/sr_cross_check.py
\end{lstlisting}
Точные вероятности, Монте-Карло и время работы:
\begin{lstlisting}[language=bash]
poetry run python experiments/sr/sr_analysis.py
\end{lstlisting}
Сводный эксперимент по методам (нужен \texttt{npSPC}):
\begin{lstlisting}[language=bash]
RSCRIPT=$RENV/bin/Rscript R_LIBS=$RENV/lib/R/library \
    poetry run python experiments/sr/sr_modes_experiment.py
\end{lstlisting}
Координаты графиков отчёта:
\begin{lstlisting}[language=bash]
poetry run python experiments/sr/sr_gen_figdata.py > docs/reports/sr_macros.tex
\end{lstlisting}

\paragraph{Распределение и моменты.} Формула \eqref{eq:sd} для $\SD(\SR)$
приведена в \cite[с.~200, табл.~4.3]{Chakraborti2019}; правило выбора $a$
и правило сигнала --- \cite[\S4.2.1.2, с.~196--202]{Chakraborti2019}.

% ================================================================== %
\begin{thebibliography}{9}
% ================================================================== %

\bibitem{Amin1995}
Amin, R.~W., Reynolds~Jr., M.~R., Bakir, S.~T. (1995).
Nonparametric quality control charts based on the sign statistic.
\emph{Communications in Statistics: Theory and Methods}, 24(6), 1597--1623.

\bibitem{Bakir2004}
Bakir, S.~T. (2004).
A distribution-free Shewhart quality control chart based on signed-ranks.
\emph{Quality Engineering}, 16(4), 613--623.

\bibitem{Gibbons2010}
Gibbons, J.~D., Chakraborti, S. (2010).
\emph{Nonparametric Statistical Inference} (5th ed.). Chapman \& Hall/CRC.

\bibitem{Wilcoxon1945}
Wilcoxon, F. (1945).
Individual comparisons by ranking methods.
\emph{Biometrics Bulletin}, 1(6), 80--83.

\bibitem{Liu1995}
Liu, R.~Y. (1995).
Control charts for multivariate processes.
\emph{Journal of the American Statistical Association}, 90(432), 1380--1387.

\bibitem{Chakraborti2019}
Chakraborti, S., Graham, M.~A. (2019).
\emph{Nonparametric Statistical Process Control}. Wiley.

\bibitem{Nascimento2026}
Nascimento, D.~C., Fons\^eca, A.~O., Fons\^eca, M.~B., Silva, M.~A.~F.,
Ferreira, P.~H. (2026).
Nonparametric SPC: Solving complex industrial monitoring tasks with low
computational cost. \emph{Array}, 30, 100828.
\url{https://doi.org/10.1016/j.array.2026.100828}.

\bibitem{Flores2021}
Flores, M., Fern\'andez-Casal, R., Naya, S., Tarr\'io-Saavedra, J. (2021).
qcr: Quality Control Review. \emph{The R Journal}.
\url{https://doi.org/10.32614/RJ-2021-034}.

\bibitem{Bakir1979}
Bakir, S.~T., Reynolds~Jr., M.~R. (1979).
A nonparametric procedure for process control based on within-group ranking.
\emph{Technometrics}, 21(2), 175--183.

\end{thebibliography}

\end{document}
